\chapter{Нев'язка Ґреда--Шафранова як метрика}\label{ch:g04}

Чи можна оцінити передбачене поле $\psi(R,Z)$ не лише за тим, наскільки воно близьке до EFIT, а й за тим, чи є воно \emph{рівновагою} взагалі? У цьому розділі йдеться про пункти E6--E8, E22, E23, рішення D6 (метрика $S'$), гіпотези C4 і C6 та побічну знахідку E33. Фізику рівняння Ґреда--Шафранова (GS) див. у~\cite{Manual}, розд.~3, а реконструкцію рівноваги й машинне навчання~--- там само, розд.~12.

\section{Що не так з \texorpdfstring{$L^2$}{L2}: мотивація}\label{sec:g04-motivation}

Метрика Fusion Equilibrium Challenge (FEC)~\cite{FEC} зважує чотири члени: $R^2_\psi$, похідні скаляри $\qnf$ і $\betaN$, відстань між LCFS і Consistency (пулований $R^2$ семи похідних скалярів). Жоден із них не перевіряє, чи задовольняє поле фізичне рівняння. Три незалежні свідчення, що це не косметика, зібрано в \texttt{04-novelty/our-contribution.md}:

\begin{established}{E22 (прочитано)}
Мала $L^2$-похибка не гарантує малої фізичної нев'язки. У~\cite{Ding2025} суто керовані (supervised) моделі мали відносну $L^2$-похибку 0,25\,\% за великих нев'язок GS~\cite{RegEst}. Навчене відображення не є проєкцією Гальоркіна, тож стійкості на кшталт леми Сеа тут немає.
\end{established}

\begin{established}{E23 (прочитано)}
Жоден фузійний бенчмарк не оцінює нев'язку рівняння в частинних похідних (PDE). Поіменно перевірено FEC~\cite{FEC}, TokaMark~\cite{TokaMark2026}, PDEBench і The Well (журнал верифікації, V-A4~\cite{RegVerif}). Найближчий контрприклад~--- VEQDB (arXiv:2609.23296). Проте її «проєктований резидуал» (projected residual) міряє збіжність \emph{розв'язувача}, а не передбачення моделі. Окремі статті нев'язку звітують, тому твердження звужено до \emph{бенчмарків}: «бенчмарки оцінюють виходи, статті оцінюють фізику»~\cite{RegVerif}.
\end{established}

Третє свідчення наше: $L^2$-середнє двох гілок розв'язку GS має нев'язку 1,22 проти $6{,}4\cdot10^{-8}$ у справжньої гілки, і $L^2$-метрика цього не бачить за побудовою (E4, розд.~\ref{ch:g03}). З цього виросло рішення D6~\cite{RegDec}: метрика $S'$ дорівнює метриці FEC плюс GS-член плюс член калібрування невизначеності (UQ). Умова, за якої D6 переглядається: перевірка V-A4 покаже, що хтось уже оцінює нев'язку. Поки що не показала.

\section{Перешкода: у поданні немає \texorpdfstring{$p'$ і $FF'$}{p′ і FF′}}\label{sec:g04-obstacle}

Рівняння GS (\cite{Manual}, розд.~3) має вигляд
\begin{equation}\label{eq:g04-gs}
  \GSop\psi\equiv R\frac{\pa}{\pa R}\Big(\frac1R\frac{\pa\psi}{\pa R}\Big)+\frac{\pa^2\psi}{\pa Z^2}
  =-\mu_0R^2p'(\psi)-FF'(\psi).
\end{equation}
Команда подає лише $\psi$ на сітці $65\times65$. Профілів $p'$ і $FF'$ у поданні немає, тож наївну нев'язку обчислити не можна.

Вихід підказує сам EFIT~\cite{Lao1985}. Коли $\psi$ зафіксовано, праву частину \eqref{eq:g04-gs} можна розкласти за базисом поліномів від $\psiN$, і вона \emph{лінійна} за коефіцієнтами. Тому питання «чи є $\psi$ рівновагою?» зводиться до лінійних найменших квадратів (МНК): чи існує хоч якась пара функцій потоку, з якою \eqref{eq:g04-gs} виконується? Відносну нев'язку після найкращого підбору визначено в \texttt{04-novelty/metric-design.md}:
\begin{equation}\label{eq:g04-g}
  g(\psi)=\min_{\{a_i\},\{b_j\}}
  \frac{\big\|\GSop\psi+\sum_{i=0}^{n_p-1}a_iR^2\psiN^{\,i}+\sum_{j=0}^{n_f-1}b_j\psiN^{\,j}\big\|_{\Omega}}
       {\|\GSop\psi\|_{\Omega}},\qquad n_p=n_f=4 .
\end{equation}
Тут $\Omega$~--- пікселі всередині LCFS, крім двох крайових рядів і стовпців сітки. Коефіцієнти $a_i$ поглинають $\mu_0p'$, а $b_j$~--- $FF'$. Значення $g=0$ означає, що $\psi$ є рівновагою для деяких поліноміальних профілів.

\paragraph{Чому гладкістю це не обійти.} Умова \eqref{eq:g04-g} вимагає, щоб $\GSop\psi$ на кожній магнітній поверхні була афінною функцією $R^2$ з коефіцієнтами, що залежать лише від $\psi$. Це і є умова поверхні потоку. Гладке, але хибне поле її не виконує автоматично.

\section{Де це в коді}\label{sec:g04-code}

\code{ourtry/04-novelty/gs_residual_probe.py}{38}{61}{(оператор $\GSop$ і область плазми)}

\begin{itemize}
\item \textbf{Рядки 38--44} (\texttt{delta\_star}): $\GSop\psi=\pa_R^2\psi-R^{-1}\pa_R\psi+\pa_Z^2\psi$, центральні різниці другого порядку через \texttt{np.gradient}. Це розкритий вираз з \eqref{eq:g04-gs}: $R\,\pa_R(R^{-1}\pa_R\psi)=\pa_R^2\psi-R^{-1}\pa_R\psi$. Похідна береться двічі, тому похибка сітки підсилюється як $1/h^2$ (див.~\S\ref{sec:g04-physics}).
\item \textbf{Рядки 47--58} (\texttt{plasma\_mask}): LCFS видобувається \emph{тією самою} процедурою скорера (\texttt{extract\_lcfs}), що й для решти метрики. Точки всередині контуру знаходить \texttt{matplotlib.path}. Два крайові ряди й стовпці сітки відкидаються, щоб шаблон різниць не виходив за край. Зауважимо: до 2026-09-29 docstring обіцяв «ерозію на одну комірку», хоча код ерозії не робить, а лише обрізає край сітки; docstring виправлено 2026-09-29~\cite{RegLog}. Точки біля самої LCFS входять в $\Omega$.
\item \textbf{Рядки 59--61}: $\psiax$ береться в O-точці, $\psib$~--- як середнє по кожній восьмій точці контуру (білінійна інтерполяція \texttt{\_bilin}). З них отримуємо $\psiN$.
\end{itemize}

\code{ourtry/04-novelty/gs_residual_probe.py}{72}{86}{(найкращий МНК-підбір профілів)}

\begin{itemize}
\item \textbf{Рядки 74--76}: якщо LCFS не знайдено або в області менше 50 пікселів, повертається \texttt{nan}. Це чесне «не можу оцінити», а не нуль.
\item \textbf{Рядки 77--82}: стовпці матриці МНК~--- $R^2\psiN^{\,i}$ (відповідають $p'$) і $\psiN^{\,j}$ (відповідають $FF'$), по чотири кожного типу. $\psiN$ обрізано до $[0,1]$.
\item \textbf{Рядки 83--86}: \texttt{lstsq} розв'язує $A\,c\approx-y$, де $y=\GSop\psi$. Результат~--- відносна норма залишку $r=y+Ac$, тобто формула \eqref{eq:g04-g}.
\end{itemize}

\code{ourtry/04-novelty/gs_residual_probe.py}{106}{122}{(драбина збурень у \texttt{main})}

Рядки 106--115: для 12 рівномірно взятих кадрів розряду DIII-D 203702 обчислюємо $g$ на істині (EFIT) і на полях з білим шумом 1\,\% і 3\,\% від $\mathrm{std}(\psi)$. Зерно генератора дорівнює номеру кадру, тож прогін відтворюваний. Рядки 117--119 додають два контрольні варіанти: гаусове згладжування ($\sigma=1{,}5$ комірки) як «правдоподібне, але хибне» поле і масштаб $\psi\times1{,}05$ як вбудовану перевірку інваріантності.

\section{Що вийшло: E6--E8}\label{sec:g04-results}

\begin{established}{E6, E7, E8}
\textbf{E6.} GS-резидуал обчислюється лише з поданого $\psi$: на істині $g=0{,}009$, з 1\,\% шуму~--- 0,650 (у~70 разів більше), з 3\,\%~--- 0,921.
\textbf{E7.} Член інваріантний до масштабу: $\psi\times1{,}05$ значення не змінює.
\textbf{E8.} Член сліпий до гладкої похибки й до низькорангового усічення: гаусове згладжування дає 0,010, PCA-5~--- 0,0093 проти 0,009 на істині (було 0,016 без скрипта; уточнено в реєстрі 2026-09-29)~\cite{RegEst}.
\end{established}

Повторний запуск 2026-09-29 (\texttt{gs\_residual\_probe.py}, venv starter kit) дав такі медіани: істина 0,0092, шум 1\,\%~--- 0,6500, шум 3\,\%~--- 0,9207, згладжування 0,0095, масштаб 0,0092 (збігається з істиною до останнього знака). Числа E6, E7 і згладжування з E8 відтворились.

\emph{Розбіжність.} Рядка PCA-5 у цьому скрипті немає. Число 0,016 записано в \texttt{metric-\allowbreak design.md} і в E8, але скрипта, що його дає, у репозиторії не знайдено (ситуація та сама, що з E14, розд.~\ref{ch:g02}). Для путівника ми відновили перевірку тими самими функціями (\texttt{guide/\allowbreak figures/\allowbreak g04\_pca5\_check.py}): PCA-базис 166 скінченних кадрів розряду 203702, усічення до 5 компонент, ті самі 12 кадрів. Медіана вийшла $g=0{,}0093$, тобто ще ближче до істини. Висновок E8 («сліпий») тільки посилюється, але саме число 0,016 не відтворюється. Зареєстровано 2026-09-29 як уточнення E8: у реєстрі тепер 0,0093 (було 0,016 без скрипта), скрипт~--- \texttt{Our try/04-novelty/pca5\_check.py} (копія цього)~\cite{RegEst,RegLog}.

\section{Фізика: чому оператор другого порядку бачить те, чого не бачить \texorpdfstring{$L^2$}{L2}}\label{sec:g04-physics}

\begin{derivation}[спектральна вага нев'язки]
Нехай похибка передбачення~--- одна фур'є-гармоніка $\delta\psi=\varepsilon\,\ee^{\mathrm{i}\vect{k}\cdot\vect{x}}$ з $kR\gg1$. Тоді $\|\delta\psi\|_2\propto\varepsilon$, тоді як $\GSop\delta\psi\approx-k^2\delta\psi$ і $\|\GSop\delta\psi\|_2\propto k^2\varepsilon$. Отже, у нев'язці спектр похибки зважено множником $k^2$ (в енергії~--- $k^4$). Для білого шуму на сітці з кроком $h$ домінує $k\sim\pi/h$, тож шум 1\,\% у $\psi$ стає нев'язкою порядку одиниці (E6). Гладка похибка з $k\sim1/a$ (де $a$~--- малий радіус) має вагу в $(a/h)^2$ разів меншу. Крім того, вісім вільних коефіцієнтів МНК поглинають ту частину гладкого $\GSop\delta\psi$, яка лежить у їхній лінійній оболонці. Звідси сліпота до згладжування й усічення (E8).
\end{derivation}

\begin{derivation}[інваріантність до масштабу і знака]
Якщо $\psi$ розв'язує \eqref{eq:g04-gs} з профілями $p(\psi)$, $F(\psi)$, то $c\psi$ розв'язує його з $\tilde p(s)=c^2p(s/c)$, $\tilde F(s)=cF(s/c)$ (вправа 3 розд.~3 у~\cite{Manual}). Для самої нев'язки: $\GSop(c\psi)=c\,\GSop\psi$, $\psiN$ не змінюється, оптимальні коефіцієнти множаться на $c$, і відношення в \eqref{eq:g04-g} лишається тим самим. Отже, $g(c\psi)=g(\psi)$ за умови, що вісь і LCFS знайдено ті самі. Для $c>0$ це виконується автоматично (E7). Для $c=-1$ пошук осі й LCFS у скорері залежить від знакової конвенції машини і має запасний (fallback) шлях (розд.~\ref{ch:g06}), тож рівність треба перевірити. Інваріантність до знака важлива, бо DIII-D і MAST зберігають $\psi$ з протилежними знаками. У \texttt{metric-design.md} (перевірка 4) її виведено аналітично, але чисельно ще не перевірено.
\end{derivation}

Разом це дає чесний висновок \texttt{metric-design.md}. GS-член \emph{доповнює} $R^2_\psi$, а не замінює його: він ловить зернисті поля і середні гілок, тобто виходи, які взагалі не є рівновагами, але пропускає гладку похибку форми. У~\texttt{README.md} челенджу прямо записано, що нейронні бейзлайни дають карти потоку «глобально правильні, локально шорсткі». Саме таку відмову GS-член бачить безпосередньо.

\section{Метрика \texorpdfstring{$S'$}{S′} (D6)}\label{sec:g04-metric}

Пропозицію з \texttt{04-novelty/metric-design.md} наводимо без змін. Базова метрика FEC має вигляд
\[
  S=0{,}55\,R^2_\psi+0{,}15\,R^2_{\qnf,\betaN}+0{,}10\,(1-D_{\mathrm{LCFS}})+0{,}20\,\mathrm{Consistency},
\]
а запропонована~---
\begin{equation}\label{eq:g04-Sprime}
  S'=w_\psi R^2_\psi+w_{qb}R^2_{\qnf,\betaN}+w_L(1-D_{\mathrm{LCFS}})+w_c\,\mathrm{Consistency}
     +w_r\,\mathrm{GS}_{\mathrm{score}}+w_u\,\mathrm{Calibration},
\end{equation}
\begin{equation}\label{eq:g04-GSscore}
  \mathrm{GS}_{\mathrm{score}}=\mathrm{clip}\big(1-g/g_{\mathrm{ref}},\,0,\,1\big),\qquad g_{\mathrm{ref}}\approx0{,}10,
\end{equation}
\begin{equation}\label{eq:g04-calib}
  \mathrm{Calibration}=1-\frac{1}{|A|}\sum_{a\in A}\frac{|\mathrm{emp}(a)-a|}{\max(a,1-a)},\qquad A=\{0{,}5;\,0{,}8;\,0{,}9;\,0{,}95\}.
\end{equation}
Стартові ваги: $w_\psi=0{,}40$, $w_{qb}=0{,}10$, $w_L=0{,}10$, $w_c=0{,}15$, $w_r=0{,}15$, $w_u=0{,}10$. Ваги «потребують емпіричної перевірки, а не голосування». Значення $g_{\mathrm{ref}}$ підібрано так, щоб істина давала $\approx0{,}91$, а шум 1\,\%~--- нуль. Для надійного калібрування 12 кадрів одного розряду замало. З восьми обов'язкових самоперевірок у \texttt{metric-design.md} виконано лише одну (інваріантність до масштабу). Решта, зокрема \texttt{perfect}$\,\to1$, \texttt{zeros}$\,\to0$, геймабельність широкими смугами і поведінка на MAST, ще відкриті.

\section{C4: чи змінить \texorpdfstring{$S'$}{S′} упорядкування? Частковий результат}\label{sec:g04-c4}

\begin{conjecture}{C4}
Додавання $\mathrm{GS}_{\mathrm{score}}$ і Calibration змінить упорядкування семи наявних бейзлайнів. \emph{Спростує:} упорядкування за $S'$ збігається з упорядкуванням за $S$, тобто нові члени нічого не додають. \emph{Вартість:} близько півдня: реалізувати $S'$ і прогнати \texttt{experiments.py}~\cite{RegConj}. \emph{Стан 2026-10-01:} половину з $\mathrm{GS}_{\mathrm{score}}$ спростовано (R12, \S\ref{sec:g04-c4test}); частина про Calibration відкрита.
\end{conjecture}

Повної перевірки ще немає. Є частковий результат (N6, 2026-09), записаний у C4 з позначкою «не вердикт». Скрипт практикуму розд.~12 методички порівнює три поля: EFIT, прогноз бейзлайну PCA+Ridge і EFIT, стиснуте до 50 головних компонент тієї ж PCA.

\code{manualcode/ch11_gs_residual_baseline.py}{44}{59}{(три варіанти $\psi$ на 36 кадрах)}

Рядки 44--48 читають кожен демо-розряд DIII-D повністю і беруть прогноз бейзлайну (\texttt{bl.predict\_row}). Рядки 49--51 лишають кадри, скінченні в обох полях, і рівномірно вибирають 12 з них. Рядок 55 будує $L^2$-проєкцію на 50 компонент PCA самої моделі, тобто майже точне поле. Рядки 56--59 записують $R^2$ і $g$ для всіх трьох варіантів тією самою функцією \texttt{gs\_inconsistency}.

\begin{figure}[htbp]\centering
\includegraphics[width=0.92\textwidth]{g04_gs_residual.jpg}
\caption[GS-нев'язка на 36 кадрах DIII-D]{GS-нев'язка $g$ на 36 кадрах трьох демо-розрядів DIII-D для EFIT, прогнозу PCA+Ridge і PCA-50-проєкції EFIT. Власний розрахунок, код: \texttt{manual/code/ch11\_gs\_residual\_baseline.py} (прогін 2026-09-29), рисунок~--- \texttt{guide/figures/g04\_fig.py}. Дані: Fusion Equilibrium Challenge (Sophelio / General Atomics), CC BY 4.0.}\label{fig:g04-residual}
\end{figure}

Прогін 2026-09-29 відтворив медіани C4 до останнього знака: $g(\mathrm{EFIT})=0{,}0103$, $g(\text{прогноз})=0{,}0174$, $g(\text{PCA-50})=0{,}0170$. Медіана $R^2$ прогнозу 0,877, $1-R^2(\text{PCA-50})=2{,}26\cdot10^{-7}$; LCFS знайдено на всіх 36 кадрах. Звідси два спостереження (рис.~\ref{fig:g04-residual}):
\begin{enumerate}
\item GS-нев'язка відрізняє обидва варіанти від EFIT, але ледь розрізняє слабкий прогноз і майже точну $L^2$-проєкцію. Це узгоджується з E8 і є ризиком для C4.
\item Висновок має фізичне пояснення (\S\ref{sec:g04-physics}). PCA-50 відкидає хвіст спектра, тобто дрібномасштабні компоненти, і $\GSop$ підсилює цю мізерну $L^2$-різницю в $k^2$ разів. Похибка прогнозу велика, але здебільшого гладка, і її $g$ майже не бачить. Отже, $g$ міряє «шорсткість», а не «хибність».
\end{enumerate}
Застереження з C4: демо-розряди, можливо, входять до 40 навчальних, тож цифри~--- ілюстрація. Упорядкування бейзлайнів прогнано 2026-10-01 (\S\ref{sec:g04-c4test}). На перших кадрах кожного розряду $R^2$ прогнозу від'ємний ($\approx-2{,}5$): це фаза наростання струму, де бейзлайн гірший за сталу.

\subsection{Перевірка C4 (2026-10-01): половина з \texorpdfstring{$\mathrm{GS}_{\mathrm{score}}$}{GS\_score}}\label{sec:g04-c4test}

Calibration потребує смуг невизначеності, яких жоден бейзлайн не видає, тож за згодою з
користувачем перевірили лише $\mathrm{GS}_{\mathrm{score}}$; Calibration піде разом із C6. Критерії
записали до запуску (\texttt{c4/THEORY.md}): \emph{стійка перестановка} --- пара моделей міняється
місцями за $S$ і за $S'_{\mathrm{GS}}$ щонайменше в 95\,\% бутстреп-повторень по тестових розрядах.
\begin{itemize}
\item $g_{\mathrm{ref}}$ відкалібровано на 28 розрядах: $g(\text{істина}+1\,\%\ \text{шуму})=0{,}634$,
  $g(\text{істина})=0{,}011$.
\item Член кадру --- \emph{відносно істини}, щоб ідеальний прогноз давав 1:
  $\mathrm{GS}_{\mathrm{score}}=\mathrm{clip}\bigl(1-(g_{\mathrm{pred}}-g_{\mathrm{true}})/(g_{\mathrm{ref}}-g_{\mathrm{true}}),0,1\bigr)$.
\item Сім бейзлайнів \texttt{experiments.py} і PCA+Ridge на розбитті C2 (навчання 0--35, тест 60--67);
  мережі --- 3 зерна. Головна вага $w_r=0{,}15/0{,}90$, скан від 0 до 0,30.
\end{itemize}

\code{ourtry/04-novelty/c4/evaluate.py}{86}{96}{(член $\mathrm{GS}_{\mathrm{score}}$ для кожного кадру)}

\noindent\emph{Рядки 91--92:} офіційні члени $S$ (через функції скорера) і $g$ кожного кадру
прогнозу. \emph{Рядки 93--94:} член відносно $g$ істини того самого кадру; де LCFS прогнозу не
знайдено, $g$ не визначено і член дорівнює нулю.

\begin{center}\small
\begin{tabular}{@{}lccccc@{}}
\toprule
модель & $S$ & $R^2_\psi$ & Consistency & медіана $g$ & $\mathrm{GS}_{\mathrm{score}}$ \\
\midrule
Linear / Ridge / PCA+Ridge & 0,19 & 0,09 & 0,27 & 0,034 & 0,95 \\
Random Forest & 0,17 & 0,11 & 0,11 & 0,017 & 0,98 \\
MLP (sklearn) & 0,64 & 0,96 & 0,12 & 0,54 & 0,20 \\
Simple MLP & 0,63 & 0,96 & 0,08 & 0,96 & 0 \\
Conv Decoder & 0,65 & 0,97 & 0,15 & 0,81 & 0 \\
UNet\_Lite & 0,66 & 0,98 & 0,18 & 0,86 & 0 \\
\bottomrule
\end{tabular}
\end{center}

\begin{figure}[htbp]\centering
\includegraphics[width=\textwidth]{g04_c4_gs.jpg}
\caption{Ліворуч: $\mathrm{GS}_{\mathrm{score}}$ проти $R^2_\psi$ для восьми моделей. Праворуч: $S'_{\mathrm{GS}}$ залежно
від ваги $w_r$ (пунктир --- головна вага). Власний розрахунок, код: \texttt{Our try/04-novelty/c4/\{evaluate,summarise\}.py}.}
\label{fig:g04-c4}
\end{figure}

За $S'_{\mathrm{GS}}$ MLP (sklearn) піднімається на перше місце ($\tau$ Кендалла $=0{,}857$), але
лише в 74--75\,\% повторень; навіть при $w_r=0{,}30$ --- 94\,\% і 84\,\%. Групи «лінійні» ($S\approx0{,}19$)
і «мережі» ($S\approx0{,}65$) розділені надто сильно, щоб член вагою до 0,30 їх переставив
(рис.~\ref{fig:g04-c4}, праворуч).

\begin{refuted}{R12 (половина гіпотези C4)}
«$\mathrm{GS}_{\mathrm{score}}$ змінить упорядкування бейзлайнів». Точковий порядок змінюється, але
жодна перестановка не тримається в 95\,\% повторень. Частина C4 про Calibration лишається відкритою
\cite{RegEst}.
\end{refuted}

\begin{established}{E38}
Карти потоку мереж із виходом-картою мають $R^2_\psi=0{,}96\text{--}0{,}98$, але $g=0{,}81\text{--}0{,}96$ ---
більше, ніж в істини з 1\,\% шуму: за рівнянням ГШ це взагалі не рівноваги. Моделі на коефіцієнтах
PCA мають $g=0{,}02\text{--}0{,}03$ завдяки гладкості, а не фізиці ($R^2_\psi\approx0{,}09$). Член
антикорелює з $R^2_\psi$ (Спірмен $-0{,}60$) \cite{RegEst}.
\end{established}

\paragraph{Що з цього випливає.} У теперішньому вигляді $\mathrm{GS}_{\mathrm{score}}$ --- не оцінка якості,
а \emph{шлюз}: він надійно відділяє «не рівновагу» від рівноваги, але нагороджує гладкість. Природніше
множити ($S\times$допуск) або ставити поріг на $g$, ніж додавати. Ризик, про який попереджала E8,
справдився.

\section{C6: калібрування невизначеності поля}\label{sec:g04-c6}

\begin{conjecture}{C6}
Одночасні конформні смуги на полі $65\times65$~--- правильна форма UQ-члена. \emph{Спростує:} емпіричне покриття виявиться недосяжним без абсурдної ширини смуг. \emph{Вартість:} близько дня~\cite{RegConj}.
\end{conjecture}

\begin{outofcorpus}[конформні передбачення]
Роздільне (split) конформне передбачення: на калібрувальній вибірці з $n$ кадрів рахуємо оцінки $s_i$. Для поля зручна max-статистика $s_i=\max_{x}|\hat\psi_i(x)-\psi_i(x)|/\hat\sigma(x)$. Беремо $\lceil(n+1)(1-\alpha)\rceil$-ту в порядку зростання оцінку $\hat s$. Смуга $\hat\psi(x)\pm\hat s\,\hat\sigma(x)$ покриває \emph{все поле одночасно} з імовірністю не менше $1-\alpha$, якщо дані обмінювані (exchangeability), причому без жодних припущень про розподіл і незалежно від якості базової моделі. Поточкові інтервали дали б 4\,225 окремих маргінальних гарантій замість однієї спільної.
\end{outofcorpus}

Отже, питання C6 справді математичне: який спосіб обрати~--- max-статистику, проєкцію на базис PCA чи конформізацію похідних скалярів замість пікселів (\texttt{metric-design.md}). Захист від геймінгу обов'язковий: нескінченно широка смуга має ідеальне покриття. Тому \eqref{eq:g04-calib} треба доповнити штрафом за ширину (нормування ширини, CRPS або логарифмічна передбачувальна щільність) \emph{до} публікації правил. Стан ніші за \texttt{our-contribution.md}: «підтверджено тонкою». GP-томографія дає UQ у замкненій формі, але відкидає баланс сил, а EFIT-Prime не доводить частотного калібрування.

\section{Побічна знахідка: функція Гріна з неправильним аргументом (E33)}\label{sec:g04-e33}

\begin{established}{E33 --- виправлено 2026-09-29}
У \texttt{Tokamak-GS-solver} функції \texttt{ellipk}/\texttt{ellipe} отримували модуль $k$ замість параметра $m=k^2$. Похибка старої $G$ сягала $+9\,\%$ біля нитки, $+42\,\%$ при $k^2=0{,}95$ ($3{,}355\cdot10^{-7}$ проти $2{,}355\cdot10^{-7}$ Вб/рад/А) і $\times8$ при $k^2=0{,}5$. Після виправлення відхилення від квадратури Біо--Савара становить $\sim10^{-15}$~\cite{RegEst,g04GSsolverChanges}.
\end{established}

\code{gs/src/numerical/compute.py}{12}{31}{(функція Гріна кругової нитки, після виправлення)}

Рядки 26--29: параметр $k^2$ обчислюється явно і передається в \texttt{special.ellipk}/\texttt{ellipe}, бо SciPy очікує саме $m=k^2$. Рядок 30 містить формулу $G=\frac{\mu_0}{2\pi}\frac{\sqrt{RR_0}}{k}\big[(2-k^2)K(k^2)-2E(k^2)\big]$. Чому помилка така велика далеко від нитки? За $k\to0$ дужка $(2-k^2)K-2E\propto k^4$ є результатом тонкого взаємного скорочення великих доданків. Якщо обчислити $K$ і $E$ у точці $k$ замість $k^2$, скорочення руйнується, і потік від далеких котушок помиляється на порядки (див.~\cite{g04GSsolverChanges}, §2). Це той самий урок, що й у \S\ref{sec:g04-physics}. Величини, отримані з \emph{різниць} великих чисел (скорочення в $G$, друга похідна в $\GSop$), підсилюють малі похибки входу. Тому для них потрібні незалежні перевірки: квадратура для $G$ (\texttt{tests/test\_green.py}) і нев'язка $g$ для $\psi$. Усі \texttt{result/PINN\_*.png}, зроблені до 29.09, використовували хибну $G$.

\begin{practicum}[C4 і C6: студентські задачі]
\textbf{Відтворити} (з теки \texttt{fusion equilibrium challenge/starter}, $\approx$20~с і $\approx$4~с):
\texttt{.venv/bin/python "../../Our try/04-novelty/gs\_residual\_probe.py"}; \texttt{.venv/bin/python ../../manual/code/ch11\_gs\_residual\_baseline.py}.

\textbf{Задача~1 (R12 $\to$ шлюз, пів дня).} Половину C4 спростовано (\S\ref{sec:g04-c4test}). Спробуйте $\mathrm{GS}_{\mathrm{score}}$ як шлюз: $S\cdot\mathbb 1[g<g_{\mathrm{ref}}]$ або м'якшу шкалу $\log g$. Чи з'являться стійкі перестановки, і чи вони розумні (лінійні моделі не мають вигравати лише завдяки гладкості)? Код і прогнози --- \texttt{Our try/04-novelty/c4/}.

\textbf{Задача~2 (самоперевірки).} Перевірте чисельно $g(-\psi)=g(\psi)$ (перевірка 4) і стійкість до порядку базису: $n_p,n_f=2\ldots8$ (перевірка 6). Як змінюється відношення $g(\text{шум }1\,\%)/g(\text{істина})$?

\textbf{Задача~3 (C6, день).} Розбийте кадри на навчальні й калібрувальні \emph{за розрядами}, а не за кадрами, інакше обмінюваність порушено. Побудуйте max-статистичну смугу для $\alpha=0{,}1$ і виміряйте емпіричне покриття та середню ширину на відкладених розрядах. Якщо покриття 90\,\% досягається лише зі смугою, ширшою за $\mathrm{std}(\psi)$, це аргумент на користь спростування C6.
\end{practicum}
